\documentclass{ximera}

\usepackage{fullpage}
\usepackage{amsthm}
\usepackage{amsmath}
\usepackage{amsfonts}
\usepackage{amssymb}
\usepackage{graphicx}
\usepackage{amsmath, amssymb}
\usepackage{lipsum}
\usepackage{enumitem}
\usepackage{hyperref}

\title{Project 3: Brandon Clark}
\author{Brandon Clark}

\begin{document}
\begin{abstract}
    Project 3 on kaleidoscope completed by Brandon Clark .
\end{abstract}

\maketitle

\section*{Overview}
The goal of this project is to use the power of matrices and linear transformations to build a virtual kaleidoscope.  



\section*{Virtual Kaleidoscopes}
\begin{center}
\includegraphics[width = 2in]{start.jpg} 
\includegraphics[width = 2in]{finish.jpg}
\end{center}

A physical kaleidoscope\footnote{There are many different kinds of kaleidoscopes.  This is just one example.} starts with a pie-shaped image like the one on the left above.  Two mirrors are held perpendicular to the paper with the image.  One mirror is held above the horizontal edge of the image and the other is held above the angled edge of the image.  The two edges and, thus, the two mirrors meet in a $60^{\circ}$ angle.  The horizontal mirror alone would produce the reflection on the left below.  The diagonal mirror alone would produce the reflection on the right below. 

\begin{center}
\includegraphics[width = 2in]{horizontal-reflection.jpg}
\includegraphics[width = 2in]{diagonal-reflection.jpg}
\end{center}

\begin{center}
 \includegraphics[width = 4in]{finish.jpg}
  \end{center}
Kaleidoscopic images like the one shown above are formed by reflections, reflections of reflections, reflections of reflections of reflections, and so on.  Each reflection produces a new pie-shaped piece (or sector) of the final image.  Why do you think there are only six pie-shaped pieces or sectors?

We will use two-dimensional vectors $\vec x = \langle{x_1, x_2}\rangle$ to denote the coordinates of each point in the original image.  Mirrors flip or reflect images.  The same point whose original coordinates were denoted by the vector $\vec x = \langle{x_1, x_2}\rangle$ will now have coordinates denoted by a new vector $\vec y = \langle{y_1, y_2}\rangle$ after reflection.  
The new coordinates are determined by a linear transformation
\begin{eqnarray*}
y_1 &=& a_{11} x_1 + a_{12} x_2\\
y_2 &=& a_{21} x_1 + a_{22} x_2
\end{eqnarray*}

$$\left[{\begin{array}{c} y_1 \\ y_2 \\ \end{array}}\right] = \left[{\begin{array}{cc} a_{11} & a_{12} \\ a_{21} & a_{22} \\ \end{array}}\right] \left[{\begin{array}{c} x_1 \\ x_2 \\ \end{array}}\right]
\quad {\rm or} \quad \vec y = A \vec x.$$

\begin{enumerate}
    \item Find the linear transformation that reflects an image by a mirror on the $x$-axis.  Express it as a matrix.  Write the matrix below.
\[
\begin{bmatrix}
1 & 0 \\
0 & -1 \\
\end{bmatrix}
\]

\item Find the linear transformation that rotates an image counterclockwise through an angle of $60^{\circ}$.  Express it as a matrix.  Write the matrix below.

\[
\begin{bmatrix}
cos(60) & -sin(60) \\
sin(60) & cos(60) \\
\end{bmatrix}=
\begin{bmatrix}
\frac{1}{2} & -\frac{\sqrt{3}}{2} \\
\frac{\sqrt{3}}{2} & \frac{1}{2} \\
\end{bmatrix}
\]

\item Find the linear transformation that reflects an image by a mirror on the diagonal line that makes an angle of $60^{\circ}$ with the $x$-axis.  Express it as a matrix.  Write the matrix below.

60 degrees reflect = 120 degrees
\[
\begin{bmatrix}
cos(120) & sin(120) \\
sin(120) & -cos(120) \\
\end{bmatrix}=
\begin{bmatrix}
-\frac{1}{2} & \frac{\sqrt{3}}{2} \\
\frac{\sqrt{3}}{2} & \frac{1}{2} \\
\end{bmatrix}
\]

\noindent Explain how you developed this matrix.
I developed this matrix by looking at the unit circle and thinking about how reflections related to angles. I identified that the reflection is ideal between 60 and 120 degrees. The negative sign shows up for cosine because a reflection flips the part of the vector that is perpendicular to the line.

\end{enumerate}

\end{document}  